\documentclass[10pt,a4paper]{amsart}
\usepackage[margin=19mm]{geometry}
\usepackage{amsmath,amssymb,mathtools}
\usepackage[scr=rsfs,cal=euler]{mathalfa}
\usepackage{xfrac}
\usepackage[unicode,hidelinks,bookmarks=false]{hyperref}
\newcommand{\ie}{i.e.\ }
\newcommand{\cf}{cf.\ }
\newcommand{\resp}{respectively\ }
\newcommand{\Subsection}{\subsection}
\providecommand{\nobreakdash}{\nobreak\hbox{-}}
\setlength{\emergencystretch}{2em}
\multlinegap0pt
\usepackage{csquotes}
\usepackage{xcolor}
\usepackage[shortlabels]{enumitem}
\usepackage{amssymb}
\usepackage{siunitx}
\usepackage{bm}
\usepackage[normalem]{ulem}
\newtheorem{lemma}{Lemma}
\newcommand*{\ldblbrace}{\{\mskip-5mu \left \{}
\newcommand*{\rdblbrace}{ \}\mskip-5mu \right \}}
\newcommand{\mean}[1]{\left\ldblbrace #1 \right\rdblbrace}
\renewcommand{\vec}[1]{\boldsymbol{#1}}
\title{On Affordable High-Order Entropy-Conservative/Stable and Well-Balanced Methods for Nonconservative Hyperbolic Systems}
\author{Marco Artiano, Hendrik Ranocha}
\date{Source: arXiv:2603.18978v3 (CC BY 4.0)}
\begin{document}
\maketitle

\noindent\emph{Source and licence.} On Affordable High-Order Entropy-Conservative/Stable and Well-Balanced Methods for Nonconservative Hyperbolic Systems. Version arXiv:2603.18978v3 (CC BY 4.0), distributed by arXiv under the Creative Commons Attribution 4.0 International licence, http://creativecommons.org/licenses/by/4.0/, as recorded in the arXiv licence metadata for this exact version and shown on the abstract page https://arxiv.org/abs/2603.18978v3. Full licence notice: \texttt{licenses/arxiv-2603.18978-CC-BY-4.0.txt}.\par\smallskip

\noindent\textit{Editorial scope.} Counted unit: the article's lemma algebraic\_lemma (source0.tex lines 330--351). The article's complete original proof of it is delivered with it (source0.tex lines 352--414, one \texttt{proof} environment of 63 lines). No mathematics is written, repaired or re-derived: the statement and the proof are the article's own lines, in the article's own notation, and no retained line is substituted or re-flowed.  No further source-local context is needed: every cross-reference of every retained line resolves inside the two delivered spans.  Nothing was omitted from any retained span, so the package carries no omission record.  No retained line contains a citation, so the wrapper carries no bibliography and no bibliography entry.  No retained line contains a figure environment, an image-inclusion command or drawing code, so no image is delivered and the package declares no asset dependency.  Editorial formatting only, in addition: an AMS \texttt{amsart} preamble that restates the article's own declarations and macros used by the retained lines, namely the environments \texttt{lemma} on separate counters, together with the standard packages those lines need.  The retained environments print in this document as Lemma 1 in order of appearance, whereas the article prints the same items with the numbering of its own full document.  The complete native source of the article as distributed in the arXiv e-print for this exact version is bundled unchanged as \texttt{sources/196744/source0.tex}.
\begin{lemma}{\label{algebraic_lemma}}
    Let $Y \subset \mathbb{R}^n$ be open, and $\vec{\omega}\colon Y \to \mathbb{R}^n$, $\vec{f}\colon Y \to \mathbb{R}^n$, $F\colon Y \to \mathbb{R}$, and $\vec{f}^{\mathrm{num}}\colon Y \times Y \to \mathbb{R}^n$ satisfying
    \begin{equation}
    \forall \vec{u} \in Y\colon \quad \vec{f}^\mathrm{num}(\vec{u},\vec{u}) = \vec{f}(\vec{u})
    \end{equation}
    be given. Then,
    $\exists F^\mathrm{num}\colon Y \times Y \to \mathbb{R}$ satisfying $\forall \vec{u} \in Y\colon F^\mathrm{num}(\vec{u},\vec{u}) = F(\vec{u})$ and
    \begin{equation}
    \forall \vec{u}_-, \vec{u}_0, \vec{u}_+ \in Y\colon \quad \vec{\omega}(\vec{u}_0) \cdot \Bigl ( \vec{f}^\mathrm{num}(\vec{u}_0, \vec{u}_+) - \vec{f}^\mathrm{num}(\vec{u}_-, \vec{u}_0) \Bigr ) \geq F^\mathrm{num}(\vec{u}_0, \vec{u}_+) - F^\mathrm{num}(\vec{u}_-, \vec{u}_0)
    \label{sufficient_condition}
    \end{equation}
    if and only if
    \begin{equation}
    \forall \vec{u}_+, \vec{u}_- \in Y\colon \quad \left (\vec{\omega}_+ - \vec{\omega}_- \right ) \cdot \vec{f}^\mathrm{num}(\vec{u}_-, \vec{u}_+) \leq \left (\vec{\omega}_+ \cdot \vec{f}_+ - F_+ \right ) - \left (\vec{\omega}_- \cdot \vec{f}_- - F_- \right ).
    \label{necessary_condition}
    \end{equation}
    Moreover, if one of the conditions is satisfied with equality, then equality holds in the other condition as well, and $F^\mathrm{num}$ in~\eqref{sufficient_condition} is determined uniquely as
    \begin{equation}
    F^{\mathrm{num}} = \mean{F} + \mean{\vec{\omega}} \cdot \vec{f}^{\mathrm{num}}  -\mean{\vec{\omega} \cdot \vec{f}}.
    \label{F_num}
    \end{equation}
\end{lemma}

\begin{proof}
We first show that~\eqref{sufficient_condition} implies~\eqref{necessary_condition}.
First, we choose  $\vec{u}_0 = \vec{u}_-$ in~\eqref{sufficient_condition}, resulting in
\begin{equation}
\vec{\omega}_- \cdot \Bigl (\vec{f}^{\mathrm{num}}(\vec{u}_-, \vec{u}_+) - \vec{f}_- \Bigr ) \geq F^{\mathrm{num}}(\vec{u}_-, \vec{u}_+) - F_-.
\label{eqminus}
\end{equation}
Choosing $\vec{u}_0 = \vec{u}_+$ in~\eqref{sufficient_condition} yields
\begin{equation}
\vec{\omega}_+ \cdot \Bigl (\vec{f}_+ - \vec{f}^{\mathrm{num}}(\vec{u}_-, \vec{u}_+) \Bigr ) \geq F_+ - F^{\mathrm{num}}(\vec{u}_-, \vec{u}_+).
\label{eqplus}
\end{equation}
Adding the two inequalities, we obtain
\begin{equation}
\left ( \vec{\omega}_- - \vec{\omega}_+ \right ) \cdot \vec{f}^{\mathrm{num}}(\vec{u}_-, \vec{u}_+) \geq \left (\vec{\omega}_- \cdot \vec{f}_- - F_- \right ) - \left (\vec{\omega}_+ \cdot \vec{f}_+ - F_+ \right ),
\end{equation}
which is equivalent to~\eqref{necessary_condition}.

If equality holds in~\eqref{sufficient_condition}, it holds in~\eqref{necessary_condition} as well.
In this case, we can subtract~\eqref{eqplus} from~\eqref{eqminus} (both with equality) to obtain
\begin{equation}
F^{\mathrm{num}}(\vec{u}_-, \vec{u}_+) = \frac{\vec{\omega}_+ + \vec{\omega}_-}{2} \cdot \vec{f}^{\mathrm{num}}(\vec{u}_-, \vec{u}_+) - \frac{\vec{\omega}_+ \cdot \vec{f}_+ + \vec{\omega}_- \cdot \vec{f}_-}{2} + \frac{F_+ + F_-}{2},
\end{equation}
i.e., $F^\mathrm{num}$ is determined uniquely by~\eqref{F_num}.

Finally, we show that~\eqref{necessary_condition} implies~\eqref{sufficient_condition}.
Abbreviating $\psi = \vec{\omega} \cdot \vec{f} - F$, \eqref{necessary_condition} is equivalent to
\begin{equation}
-\left (\vec{\omega}_+ - \vec{\omega}_- \right ) \cdot \vec{f}^\mathrm{num}(\vec{u}_-, \vec{u}_+) \ge - \left(\psi_+ - \psi_-\right).
\label{necessary_condition_psi}
\end{equation}
Following Tadmor, we compute
\begin{equation}
\begin{aligned}
    &\qquad
    \vec{\omega}(\vec{u}_0) \cdot \Bigl ( \vec{f}^\mathrm{num}(\vec{u}_0, \vec{u}_+) - \vec{f}^\mathrm{num}(\vec{u}_-, \vec{u}_0) \Bigr )
    \\
    &=
    \frac{\vec{\omega}_0 + \vec{\omega}_+}{2} \cdot \vec{f}^\mathrm{num}(\vec{u}_0, \vec{u}_+)
    - \frac{\vec{\omega}_+ - \vec{\omega}_0}{2} \cdot \vec{f}^\mathrm{num}(\vec{u}_0, \vec{u}_+)
    \\
    &\quad
    - \frac{\vec{\omega}_- + \vec{\omega}_0}{2} \cdot \vec{f}^\mathrm{num}(\vec{u}_-, \vec{u}_0)
    - \frac{\vec{\omega}_0 - \vec{\omega}_-}{2} \cdot \vec{f}^\mathrm{num}(\vec{u}_-, \vec{u}_0)
    \\
    &\stackrel{\eqref{necessary_condition_psi}}{\ge}
    \frac{\vec{\omega}_0 + \vec{\omega}_+}{2} \cdot \vec{f}^\mathrm{num}(\vec{u}_0, \vec{u}_+)
    - \frac{\psi_+ - \psi_0}{2}
    - \frac{\vec{\omega}_- + \vec{\omega}_0}{2} \cdot \vec{f}^\mathrm{num}(\vec{u}_-, \vec{u}_0)
    - \frac{\psi_0 - \psi_-}{2}
    \\
    &=
    \frac{\vec{\omega}_0 + \vec{\omega}_+}{2} \cdot \vec{f}^\mathrm{num}(\vec{u}_0, \vec{u}_+)
    - \frac{\psi_+ + \psi_0}{2}
    - \frac{\vec{\omega}_- + \vec{\omega}_0}{2} \cdot \vec{f}^\mathrm{num}(\vec{u}_-, \vec{u}_0)
    + \frac{\psi_0 + \psi_-}{2}
    \\
    &=
    F^{\mathrm{num}}(\vec{u}_0, \vec{u}_+) - F^{\mathrm{num}}(\vec{u}_-, \vec{u}_0),
\end{aligned}
\end{equation}
where we have used the definition of $F^{\mathrm{num}}$ in~\eqref{F_num} in the last step.
\end{proof}

\end{document}
